\begin{lemma}
Let $\nu$ be a finite measure on a measurable space $\Omega$, let $s$ be a
finite set of indices, and let $P_i$ be predicates on $\Omega$ whose events
$A_i=\{\omega:P_i(\omega)\}$ are measurable for $i\in s$.
The real-valued count $C(\omega)=|\{i\in s:P_i(\omega)\}|$ is integrable and
\[
\int C\,d\nu=\sum_{i\in s}\nu(A_i),
\]
where the measures on the right are interpreted as real numbers.
\end{lemma}
\begin{proof}
Pointwise, $C=\sum_{i\in s}\mathbf1_{A_i}$. The constant function one is
integrable since $\nu$ is finite. Its restriction by the indicator of each
measurable $A_i$ is integrable, so the finite sum is integrable as well.
Linearity of the integral for a finite sum of integrable functions gives
$\int C\,d\nu=\sum_{i\in s}\int\mathbf1_{A_i}\,d\nu$.
Each summand is the real measure of $A_i$, proving the formula, also when
$s$ is empty.
\end{proof}
